% =============================================================================
% TP1 --- Société Orion : conception et implémentation d'un Data Warehouse
% Compilation : tectonic rapport.tex
% =============================================================================
\documentclass[11pt,a4paper]{article}

% --- Encodage / langue -------------------------------------------------------
\usepackage[T1]{fontenc}
\usepackage[utf8]{inputenc}
\usepackage[french]{babel}
\usepackage{lmodern}
\usepackage{microtype}
\usepackage{amsmath}
\usepackage{amssymb}
\usepackage{adjustbox}

% --- Mise en page ------------------------------------------------------------
\usepackage[a4paper,margin=2.2cm]{geometry}
\usepackage{parskip}
\usepackage{enumitem}
\setlist{nosep,topsep=2pt}
\usepackage{array}
\usepackage{booktabs}
\usepackage{longtable}
\usepackage{tabularx}
\usepackage{caption}
\captionsetup{font=small,labelfont=bf}

% --- Couleurs ----------------------------------------------------------------
\usepackage{xcolor}
\definecolor{cBlue}{HTML}{1F4E79}
\definecolor{cTeal}{HTML}{2E8B8B}
\definecolor{cOrange}{HTML}{C8702A}
\definecolor{cGreen}{HTML}{2E7D32}
\definecolor{cGrey}{HTML}{555555}
\definecolor{cBgFact}{HTML}{FFE7CE}
\definecolor{cBgDim}{HTML}{D6E8F5}
\definecolor{cBgOps}{HTML}{E8F0E8}
\definecolor{cBgInfra}{HTML}{F5EAD6}

% --- Code listings -----------------------------------------------------------
\usepackage{listings}
\lstdefinelanguage{dockerfile}{
  morekeywords={FROM,RUN,CMD,EXPOSE,ENV,ADD,COPY,ENTRYPOINT,VOLUME,USER,WORKDIR,ARG,ONBUILD,STOPSIGNAL,LABEL,SHELL,HEALTHCHECK},
  morecomment=[l]{\#}, sensitive=true
}
\lstdefinelanguage{yaml}{
  keywords={services,image,container_name,environment,ports,volumes,depends_on,
            healthcheck,networks,profiles,build,restart,driver,version,
            condition,test,interval,retries},
  morecomment=[l]{\#}, sensitive=true,
  morestring=[b]', morestring=[b]"
}
\lstdefinestyle{base}{
  basicstyle=\ttfamily\footnotesize,
  keywordstyle=\color{cBlue}\bfseries,
  commentstyle=\color{cGreen}\itshape,
  stringstyle=\color{cOrange},
  numberstyle=\tiny\color{cGrey},
  numbers=left, numbersep=6pt,
  frame=single, rulecolor=\color{black!30},
  framesep=4pt, breaklines=true, breakatwhitespace=true,
  showstringspaces=false, tabsize=2, columns=fullflexible, keepspaces=true,
  belowskip=4pt, aboveskip=4pt,
}
\lstdefinestyle{sql}{style=base,language=SQL}
\lstdefinestyle{py}{style=base,language=Python}
\lstdefinestyle{yml}{style=base,language=yaml}
\lstdefinestyle{dock}{style=base,language=dockerfile}
\lstdefinestyle{shell}{style=base,language=bash}

% --- Graphiques (diagrammes Mermaid pré-rendus) -------------------------------
\usepackage{graphicx}
\graphicspath{{build/}{screenshots/}}
\newcommand{\diag}[1]{%
  \begin{adjustbox}{max width=\textwidth,max totalheight=0.78\textheight,center}%
    \includegraphics{#1}%
  \end{adjustbox}%
}

% --- Placeholder pour captures d'écran ---------------------------------------
% Utilisation : \screenshot{fichier.png}{Légende de la capture}
%   - si screenshots/fichier.png existe -> il est inclus
%   - sinon -> un cadre vide « À insérer » est affiché à sa place
\usepackage{tikz}
\usepackage{ifthen}

\newcommand{\screenshot}[3][0.85]{%
  \begin{figure}[ht]
    \centering
    \IfFileExists{screenshots/#2}{%
      \includegraphics[width=#1\textwidth,
                       max totalheight=0.6\textheight,
                       keepaspectratio]{#2}%
    }{%
      \fbox{\begin{minipage}[c][6cm][c]{#1\textwidth}
        \centering
        {\large\bfseries\color{cBlue}\textbullet\ Capture d'écran à insérer}\\[6pt]
        {\small\ttfamily screenshots/#2}\\[10pt]
        \emph{#3}
      \end{minipage}}%
    }
    \caption{#3}
    \label{fig:scr:#2}
  \end{figure}
}

% --- Hyperref (en dernier) ---------------------------------------------------
\usepackage[colorlinks=true, linkcolor=cBlue, urlcolor=cTeal,
            citecolor=cTeal, breaklinks=true]{hyperref}

% --- Macros ------------------------------------------------------------------
\newcommand{\code}[1]{\texttt{\small #1}}
\newcommand{\file}[1]{\texttt{\small\itshape #1}}
\newcommand{\Q}[1]{\par\smallskip\noindent\textbf{Question \arabic{question}.}~#1\par\smallskip}
\newcounter{question}
\newcommand{\question}[1]{\refstepcounter{question}\subsection*{Q\arabic{question}. #1}\addcontentsline{toc}{subsection}{Q\arabic{question}. #1}}

\emergencystretch=3em
\sloppy

% =============================================================================
\begin{document}
% =============================================================================

\begin{titlepage}
  \centering
  \vspace*{2cm}
  {\Large\sffamily ENEAM --- Master Entrepôts de Données}\par
  \vspace{0.6cm}
  \rule{0.7\textwidth}{0.6pt}\par
  \vspace{1cm}
  {\Huge\bfseries\sffamily TP n\textsuperscript{o} 1\par}
  \vspace{0.4cm}
  {\LARGE\sffamily Société Orion\par}
  \vspace{0.4cm}
  {\Large\itshape Conception et implémentation\\ d'un Data Warehouse opérationnel\par}
  \vspace{1cm}
  \rule{0.7\textwidth}{0.6pt}\par
  \vspace{1.5cm}
  \begin{minipage}{0.85\textwidth}\centering
    \small
    Modélisation opérationnelle 3NF, transition vers la modélisation
    dimensionnelle en étoile, génération de données fictives via Faker,
    ETL implémenté en \emph{procédures stockées T-SQL sur SQL Server} et
    orchestré depuis Python, déploiement Docker complet (Postgres OLTP +
    Postgres DWH + SQL Server), réponses analytiques aux 15 questions du
    cahier des charges.
  \end{minipage}
  \vfill
  {\sffamily \today}
\end{titlepage}

\tableofcontents
\newpage

% =============================================================================
\section{Contexte et objectifs}
% =============================================================================

\subsection{La société Orion}

Orion est une société fictive mondiale, spécialisée dans la commercialisation
d'articles de sport et d'extérieur. Son siège social est aux États-Unis et
elle gère des filiales en Belgique, Pays-Bas, Allemagne, Royaume-Uni,
Danemark, France, Italie, Espagne et Australie.

Les ventes se font via trois canaux : magasins physiques, catalogue et
internet. Le programme de fidélité \emph{Orion Star Club} confère des
avantages aux clients qui en sont membres. L'historique disponible va du
\textbf{1\textsuperscript{er} janvier 1998 au 31 décembre 2002}.

\paragraph{Volumétrie cible.}
\begin{itemize}
  \item $\sim$\,5\,500 références produit
  \item $\sim$\,90\,000 clients (pas tous actifs)
  \item $\sim$\,980\,000 commandes
  \item $\sim$\,600 à 800 employés
  \item 64 fournisseurs (un seul fournisseur par produit)
\end{itemize}

\subsection{Objectifs décisionnels}

Les questions auxquelles le système doit répondre se regroupent en cinq
familles :
\begin{description}[style=nextline,leftmargin=1.4cm]
  \item[Performance produit] best-sellers, perte de vitesse, contribution au
    chiffre d'affaires, marges par groupe.
  \item[Effet remise] impact des remises sur les ventes et la marge.
  \item[Performance commerciale] commerciaux les plus performants par pays,
    sexe, âge, salaire.
  \item[Segmentation client] groupes de clients, clients les plus rentables.
  \item[Analyses statistiques] moyennes et écarts-types du chiffre d'affaires,
    variables explicatives, comparaison hommes/femmes parmi les commerciaux.
\end{description}

\subsection{Démarche retenue}

La construction de l'entrepôt a été décomposée en sept étapes :
\begin{enumerate}
  \item Modélisation \textbf{opérationnelle} (3NF) du système source
        --- section~\ref{sec:oltp}.
  \item \textbf{Transition} OLTP $\to$ modélisation dimensionnelle (méthode
        Kimball en 4 étapes) --- section~\ref{sec:transition}.
  \item Modélisation \textbf{dimensionnelle} (étoile) de l'entrepôt
        --- section~\ref{sec:dwh}.
  \item Provisionnement de l'environnement via \textbf{Docker Compose}
        (Postgres OLTP + Postgres DWH + SQL Server)
        --- section~\ref{sec:docker}.
  \item Génération de jeux de données réalistes via \textbf{Faker}
        --- section~\ref{sec:gen}.
  \item \textbf{ETL T-SQL sur SQL Server} et son orchestration Python
        --- section~\ref{sec:etl}.
  \item Rédaction des \textbf{requêtes analytiques} répondant aux questions
        métier --- section~\ref{sec:answers}.
\end{enumerate}

\newpage
% =============================================================================
\section{Modélisation opérationnelle (OLTP)}
\label{sec:oltp}
% =============================================================================

\subsection{Principes}

L'objectif du modèle opérationnel est de soutenir les opérations courantes :
prise de commande, gestion des stocks, gestion RH, suivi des fournisseurs.
Les contraintes fortes y sont l'\emph{intégrité référentielle} et la
\emph{non-redondance}. Le modèle est en troisième forme normale (3NF) et
range toutes les tables dans un schéma PostgreSQL nommé \code{ops}.

\subsection{Carte des sous-domaines}

Le système se décompose en quatre sous-domaines fortement couplés via la
fact transactionnelle (les commandes).

\begin{figure}[h]
\centering
\diag{01-domain-map}
\caption{Carte de domaines : les commandes constituent la fact transactionnelle alimentée par les trois domaines référentiels.}
\end{figure}

\subsection{Sous-domaine Géographie \& Organisation}

Deux hiérarchies parallèles cohabitent : la géographie commerciale (utilisée
par les fournisseurs et les clients) et la hiérarchie organisationnelle des
employés (5 niveaux comme imposé par l'énoncé).

\begin{figure}[h]
\centering
\diag{02-oltp-geo-org}
\caption{Sous-domaine \emph{Géographie \& Organisation}. La table \texttt{employee} hérite de la position dans les deux hiérarchies via \texttt{org\_group\_id} et \texttt{city\_id}, et porte une auto-référence \texttt{manager\_id}.}
\end{figure}

\subsection{Sous-domaine Produits \& Fournisseurs}

La hiérarchie produit comporte les 4 niveaux requis. L'historique de prix
et les remises sont externalisés dans deux tables distinctes pour conserver
le grain temporel.

\begin{figure}[h]
\centering
\diag{03-oltp-products}
\caption{Sous-domaine \emph{Produits \& Fournisseurs}. La hiérarchie comporte 4 niveaux ; \texttt{product\_price\_history} et \texttt{product\_discount} capturent les variations temporelles (prix et remises).}
\end{figure}

\subsection{Sous-domaine Clients \& Commandes}

C'est le c\oe{}ur transactionnel. La carte de fidélité \emph{Orion Star Club}
est externalisée pour ne pas polluer la table \code{client} (relation 1-1
optionnelle).

\begin{figure}[h]
\centering
\diag{04-oltp-customers-orders}
\caption{Sous-domaine \emph{Clients \& Commandes}. La table des lignes a une clé primaire composite (\emph{order\_id}, \emph{line\_no}) ; chaque ligne pointe vers un produit, et l'en-tête de commande vers un client, un commercial et un canal.}
\end{figure}

\subsection{Choix structurants}

\begin{itemize}
  \item \textbf{Hiérarchie produit} (4 niveaux) et \textbf{hiérarchie
    organisationnelle} (5 niveaux) sont \emph{normalisées} : une table par
    niveau, FK \og enfant vers parent \fg{}. Cela évite la duplication des
    libellés et simplifie les renommages.
  \item \textbf{Historique de prix} dans \code{product\_price\_history} :
    on retrouve le prix appliqué au moment de la commande en sélectionnant
    la fenêtre \texttt{(start\_date, end\_date)} contenant
    \texttt{order\_date}. La fenêtre courante a \texttt{end\_date IS NULL}.
  \item \textbf{Remises} : \code{product\_discount} avec validité bornée,
    permettant les promotions saisonnières.
  \item \textbf{Auto-référence} \code{employee.manager\_id} pour la
    hiérarchie de management (n+1).
  \item \textbf{Carte de fidélité} : table dédiée plutôt qu'attribut booléen,
    pour stocker numéro de carte et date d'émission.
  \item Tous les montants sont en \textbf{USD} (cf. énoncé).
\end{itemize}

\subsection{Listing : extrait du DDL OLTP}

L'intégralité du schéma est dans \file{oltp/init/01\_schema.sql}. Voici la
table centrale des commandes et l'historique de prix.

\begin{lstlisting}[style=sql,caption={Tables de commandes (extrait)},label={lst:order}]
CREATE TABLE commande (
    commande_id    BIGSERIAL PRIMARY KEY,
    date_commande  DATE NOT NULL,
    client_id      BIGINT   NOT NULL REFERENCES client(client_id),
    employe_id     BIGINT   NOT NULL REFERENCES employe(employe_id),
    canal_id       SMALLINT NOT NULL REFERENCES canal_vente(canal_id),
    montant_total  NUMERIC(14,2) NOT NULL DEFAULT 0,
    cree_le        TIMESTAMP NOT NULL DEFAULT now()
);

CREATE TABLE ligne_commande (
    commande_id     BIGINT   NOT NULL REFERENCES commande(commande_id) ON DELETE CASCADE,
    numero_ligne    SMALLINT NOT NULL,
    produit_id      BIGINT   NOT NULL REFERENCES produit(produit_id),
    quantite        NUMERIC(12,3) NOT NULL CHECK (quantite > 0),
    prix_unitaire   NUMERIC(12,2) NOT NULL CHECK (prix_unitaire >= 0),
    cout_unitaire   NUMERIC(12,2) NOT NULL CHECK (cout_unitaire >= 0),
    pct_remise      NUMERIC(5,4)  NOT NULL DEFAULT 0
                    CHECK (pct_remise BETWEEN 0 AND 1),
    PRIMARY KEY (commande_id, numero_ligne)
);
\end{lstlisting}

\begin{lstlisting}[style=sql,caption={Historique de prix avec validité ouverte},label={lst:price}]
CREATE TABLE historique_prix (
    prix_id      BIGSERIAL PRIMARY KEY,
    produit_id   BIGINT NOT NULL REFERENCES produit(produit_id) ON DELETE CASCADE,
    date_debut   DATE NOT NULL,
    date_fin     DATE,
    cout         NUMERIC(12,2) NOT NULL CHECK (cout >= 0),
    prix_vente   NUMERIC(12,2) NOT NULL CHECK (prix_vente >= 0),
    CHECK (date_fin IS NULL OR date_fin >= date_debut)
);
CREATE INDEX idx_historique_prix_periode
    ON historique_prix(produit_id, date_debut, date_fin);
\end{lstlisting}

\newpage
% =============================================================================
\section{Transition OLTP $\to$ modélisation dimensionnelle}
\label{sec:transition}
% =============================================================================

Cette section formalise \emph{comment} on est passé du modèle 3NF à un schéma
en étoile. La démarche suit la \textbf{méthode Kimball en quatre étapes}
(Kimball \& Ross, \emph{The Data Warehouse Toolkit}, 3\textsuperscript{e} éd.).

\subsection{Étape 1 --- Identifier le processus métier}

Le processus métier au cœur des questions de l'énoncé (p. 3) est :
\begin{quote}
\itshape
\og Vendre un article à un client, via un canal donné, par un commercial
donné, à une date donnée. \fg{}
\end{quote}
C'est ce processus qui produit la mesure principale (chiffre d'affaires) et
qui sera matérialisé par la table de faits.

\subsection{Étape 2 --- Choisir le grain}

Choix structurant. Trois grains candidats ont été envisagés :

\begin{table}[h]
\centering\small
\begin{tabularx}{\textwidth}{l X X}\toprule
\textbf{Grain candidat} & \textbf{Pour} & \textbf{Contre}\\\midrule
Commande complète       & Volume modeste                     & Perd la dimension produit \\
\textbf{Ligne de commande \checkmark} & Conserve le détail produit         & Volume = $\sim$\,980\,000 lignes (acceptable) \\
Ligne consolidée jour   & Très petit volume                  & Perd commande, client précis, canal \\\bottomrule
\end{tabularx}
\caption{Choix du grain : la \textbf{ligne de commande} est retenue.}
\end{table}

Toutes les mesures de \code{fact\_sales} seront additives au niveau ligne.

\subsection{Étape 3 --- Identifier les dimensions}

À partir des questions analytiques, sept dimensions sont identifiées.

\begin{table}[h]
\centering\small
\begin{tabularx}{\textwidth}{X l}\toprule
\textbf{Question analytique (énoncé)}                       & \textbf{Dimension associée}\\\midrule
\og Quels produits se vendent le mieux ? \fg{}              & \code{dim\_produit}\\
\og {[\dots]} par pays et année donnés \fg{}                & \code{dim\_geographie} + \code{dim\_date}\\
\og Marge par groupe de produit \fg{}                       & \code{dim\_produit} (hiérarchie aplatie)\\
\og Commerciaux par sexe, âge, salaire, pays \fg{}          & \code{dim\_employe} (avec attributs RH)\\
\og Quels groupes de clients sont identifiés ? \fg{}        & \code{dim\_client} (groupe + démographie)\\
\og Fournisseurs proposant des produits rentables \fg{}     & \code{dim\_fournisseur}\\
\og Différence H/F entre les commerciaux \fg{}              & \code{dim\_employe.sexe}\\
Analyse par canal de vente                                  & \code{dim\_canal}\\\bottomrule
\end{tabularx}
\caption{Dérivation des dimensions à partir des questions de l'énoncé.}
\end{table}

\subsubsection*{Aplatissement des hiérarchies}

Les hiérarchies normalisées de l'OLTP sont \textbf{aplaties} dans la
dimension. Exemple produit : 4 tables OLTP $\to$ 1 dimension.

\begin{verbatim}
ops.ligne_produit ─┐
ops.categorie_produit ─┐         dim_produit (attributs aplatis)
ops.groupe_produit ─┐  ├──────► - nom_produit
ops.produit ───────┴──          - nom_groupe_produit
                                - nom_categorie_produit
                                - nom_ligne_produit
                                - nom_fournisseur
\end{verbatim}

Idem pour la géographie (4 tables $\to$ 1 dimension) et l'organisation RH
(5 tables $\to$ attributs aplatis dans \code{dim\_employee}). Cela permet :
\begin{itemize}
  \item des \textbf{drill-down} rapides en SQL (\code{GROUP BY product\_line\_name},
        \code{GROUP BY product\_group\_name}, \dots) ;
  \item une \textbf{lisibilité} accrue pour les outils BI ;
  \item un coût en redondance assumé --- c'est précisément l'objet d'un DWH.
\end{itemize}

\subsubsection*{Surrogate keys}

Chaque dimension porte une \textbf{clé technique} (surrogate key) générée à
l'insertion : \code{cle\_produit}, \code{cle\_client}, \code{cle\_employe},
\dots\ La clé naturelle OLTP est conservée comme attribut
(\code{produit\_id}, \code{client\_id}, \dots) pour la traçabilité, mais
\textbf{n'est jamais utilisée comme FK dans \code{fait\_ventes}}.

Justifications :
\begin{itemize}
  \item Découple le DWH des changements OLTP (renommage, fusion d'IDs).
  \item Indispensable pour SCD2 : plusieurs versions d'un client ont le
        \emph{même} \code{client\_id} mais des \code{cle\_client}
        \emph{différents}.
\end{itemize}

\subsection{Étape 4 --- Identifier les faits}

Les mesures sont calculées à partir de la ligne de commande :

\begin{table}[h]
\centering\small
\begin{tabular}{l l l}\toprule
\textbf{Mesure DWH} & \textbf{Source / formule}                       & \textbf{Type}\\\midrule
quantite            & \code{ligne\_commande.quantite}                 & additive\\
prix\_unitaire      & \code{ligne\_commande.prix\_unitaire}           & non-additive\\
cout\_unitaire      & \code{ligne\_commande.cout\_unitaire}           & non-additive\\
pct\_remise         & \code{ligne\_commande.pct\_remise}              & non-additive\\
montant\_brut       & quantite $\times$ prix\_unitaire                & \textbf{additive}\\
montant\_remise     & brut $\times$ pct\_remise                       & \textbf{additive}\\
montant\_net        & brut $-$ remise                                 & \textbf{additive (CA)}\\
montant\_cout       & quantite $\times$ cout\_unitaire                & \textbf{additive}\\
montant\_marge      & net $-$ cout                                    & \textbf{additive (marge)}\\\bottomrule
\end{tabular}
\caption{Mesures de \texttt{fait\_ventes} et leur additivité.}
\end{table}

\paragraph{Pourquoi pré-calculer net / margin ?}
\begin{itemize}
  \item évite les recalculs à la volée $\to$ meilleure performance ;
  \item garantit qu'analystes et rapports parlent du \emph{même} CA / de la
        \emph{même} marge.
\end{itemize}

\paragraph{Degenerate dimensions.}
\code{commande\_id} et \code{numero\_ligne} sont conservés dans
\code{fait\_ventes} sans table dimensionnelle associée --- ce sont des
\textbf{degenerate dimensions} : utiles pour drill-down jusqu'à la commande,
mais sans attribut propre.

\subsection{Bus matrix}

\begin{table}[h]
\centering\footnotesize
\begin{tabular}{l ccccccc}\toprule
\textbf{Processus métier}    & date & produit & client & employe & fournisseur & canal & geographie\\\midrule
Vente (ligne de commande)    & \checkmark & \checkmark & \checkmark & \checkmark & \checkmark & \checkmark & \checkmark (client)\\\bottomrule
\end{tabular}
\caption{Bus matrix actuel : un seul processus. L'ajout futur de \emph{Réception fournisseur} ou \emph{Mouvement de stock} partagerait dim\_produit, dim\_fournisseur, dim\_date, dim\_geographie.}
\end{table}

\subsection{Vue synthétique de la transition}

\begin{lstlisting}[style=base,basicstyle=\ttfamily\scriptsize,frame=single]
OLTP (3NF · ops)                       ETL (T-SQL · OrionETL)              DWH (étoile · dw)
─────────────────                      ─────────────────────                ─────────────────
ops.ligne_produit       ┐
ops.categorie_produit   │ aplatit     dim.dim_produit (SCD1)                dw.dim_produit
ops.groupe_produit      ├──────────►  (full reload)
ops.produit             ┘
ops.fournisseur                       dim.dim_fournisseur (SCD1)            dw.dim_fournisseur
ops.client + groupe + ville + ...     dim.dim_client (SCD2)                 dw.dim_client
ops.employe + org_* + manager         dim.dim_employe (SCD2)                dw.dim_employe
ops.continent/pays/region/ville       dim.dim_geographie (SCD1)             dw.dim_geographie
ops.canal_vente                       dim.dim_canal  (SCD1)                 dw.dim_canal
(série de dates 1997-2003)            dim.dim_date     (statique)           dw.dim_date
ops.commande + ligne_commande         fact.fait_ventes (incr. + mesures)    dw.fait_ventes
                                      etl.watermark, etl.run_log            (méta)
\end{lstlisting}

\screenshot{transition-pgadmin-vs-mssms.png}{Vue côte à côte : le schéma
\code{ops} normalisé sur pgAdmin (à gauche) et la base \code{OrionETL}
(\code{staging}, \code{dim}, \code{fact}) sur SSMS (à droite).}

\newpage
% =============================================================================
\section{Modélisation dimensionnelle (DWH)}
\label{sec:dwh}
% =============================================================================

\subsection{Étoile vs flocon}

Pour cet entrepôt, le \textbf{schéma en étoile} a été préféré au flocon pour
trois raisons :
\begin{enumerate}
  \item \textbf{Performance} de lecture : moins de jointures sur les requêtes
    OLAP, qui dominent la charge.
  \item \textbf{Lisibilité} pour les analystes : une dimension =
    une seule table avec tous les attributs hiérarchiques aplatis.
  \item \textbf{Simplicité d'ETL} : les hiérarchies ne sont reconstruites
    qu'une fois, à l'extraction.
\end{enumerate}
La sur-redondance induite est négligeable face aux gains de simplicité.

\subsection{Schéma en étoile}

\begin{figure}[h]
\centering
\diag{05-dwh-star}
\caption{Schéma en étoile centré sur \texttt{fact\_sales} (grain = ligne de commande). Les dimensions SCD2 sont identifiées explicitement.}
\end{figure}

\subsection{Stratégie SCD (Slowly Changing Dimensions)}

Toutes les dimensions ne sont pas traitées de la même manière : certaines
peuvent être écrasées (SCD1), d'autres doivent conserver l'historique des
versions (SCD2).

\begin{table}[h]
\centering\small
\begin{tabularx}{\textwidth}{l c X}
\toprule
\textbf{Dimension}     & \textbf{SCD} & \textbf{Justification}\\\midrule
dim\_date              & --   & Statique, peuplée une seule fois (1997--2003).\\
dim\_produit           & 1    & Le nom et la hiérarchie évoluent peu ; on écrase.\\
dim\_client            & 2    & Le groupe d'achat et l'adresse évoluent ; on veut analyser les ventes \emph{telles qu'à l'époque}.\\
dim\_employe           & 2    & Salaire, section, manager évoluent ; la RH demande la traçabilité.\\
dim\_fournisseur       & 1    & Très peu de changements en pratique.\\
dim\_canal             & 1    & Référentiel quasi-statique.\\
dim\_geographie        & 1    & Référentiel.\\
\bottomrule
\end{tabularx}
\caption{Choix SCD par dimension.}
\end{table}

\paragraph{Implémentation SCD2.}
Chaque dimension SCD2 porte les colonnes :
\code{effectif\_du}, \code{effectif\_au}, \code{est\_courant},
\code{hash\_ligne} (sha256 stable des attributs SCD).
Le job ETL compare le hash de la ligne source au hash de la version courante
en base. S'ils diffèrent, l'ancienne version est close
(\texttt{est\_courant = FALSE}, \texttt{effectif\_au = ajd}) et une nouvelle
version est insérée.

\subsection{Mesures de la fact}

Les mesures sont pré-calculées au chargement pour éviter des opérations
arithmétiques répétées à la lecture :

\[
\begin{array}{rcl}
  \text{montant\_brut}   &=& \text{quantite} \times \text{prix\_unitaire}\\
  \text{montant\_remise} &=& \text{montant\_brut} \times \text{pct\_remise}\\
  \text{montant\_net}    &=& \text{montant\_brut} - \text{montant\_remise}\\
  \text{montant\_cout}   &=& \text{quantite} \times \text{cout\_unitaire}\\
  \text{montant\_marge}  &=& \text{montant\_net} - \text{montant\_cout}
\end{array}
\]

\code{montant\_net} est la convention retenue pour le \emph{chiffre d'affaires}
et \code{montant\_marge} pour la \emph{marge}. Les mesures additives
(quantite, montant\_*) peuvent être agrégées librement ; les non-additives
(prix\_unitaire, cout\_unitaire, pct\_remise) ne sont utilisables qu'au grain
ligne ou via une moyenne pondérée.

\newpage
% =============================================================================
\section{Architecture technique : stack Docker}
\label{sec:docker}
% =============================================================================

\subsection{Vue d'ensemble}

L'environnement complet est provisionné par un unique
\file{docker-compose.yml} qui démarre \textbf{six services} sur un réseau
interne \code{orion-net} :

\begin{itemize}
  \item \code{postgres-oltp} : la base opérationnelle \emph{Orion} (Postgres) ;
  \item \code{postgres-dwh}  : l'entrepôt de données en étoile (Postgres) ;
  \item \code{mssql-etl}     : SQL Server 2022, qui héberge la base
        \code{OrionETL} et toutes les procédures stockées T-SQL de
        transformation ;
  \item \code{orchestrator}  : conteneur Python long-running, transporte les
        données entre Postgres et SQL Server ;
  \item \code{data-gen}      : conteneur Python one-shot peuplant l'OLTP via
        Faker (\code{--profile seed}) ;
  \item \code{pgadmin}       : UI web pour explorer les deux bases Postgres.
\end{itemize}

\begin{figure}[h]
\centering
\diag{06-docker-stack}
\caption{Stack Docker : Postgres OLTP, SQL Server hébergeant la logique ETL, Postgres DWH, orchestrateur Python qui relie les trois, plus pgAdmin et data-gen.}
\end{figure}

\screenshot{docker-compose-ps.png}{Sortie de \texttt{docker compose ps} : les six services à l'état \emph{Up (healthy)}.}

\subsection{docker-compose.yml (extrait)}

\begin{lstlisting}[style=yml,caption={Services principaux (\texttt{docker-compose.yml})},label={lst:compose}]
services:
  postgres-oltp:
    image: postgres:16-alpine
    environment:
      POSTGRES_DB: ${OLTP_DB}
      POSTGRES_USER: ${OLTP_USER}
      POSTGRES_PASSWORD: ${OLTP_PASSWORD}
    ports: ["${OLTP_PORT}:5432"]
    volumes:
      - oltp_data:/var/lib/postgresql/data
      - ./oltp/init:/docker-entrypoint-initdb.d:ro
    healthcheck:
      test: ["CMD-SHELL", "pg_isready -U ${OLTP_USER} -d ${OLTP_DB}"]

  postgres-dwh:
    image: postgres:16-alpine
    environment: { POSTGRES_DB: ${DWH_DB}, ... }
    volumes:
      - dwh_data:/var/lib/postgresql/data
      - ./dwh/init:/docker-entrypoint-initdb.d:ro

  mssql-etl:                       # moteur ETL : SQL Server + procs T-SQL
    build: ./mssql-etl
    environment:
      ACCEPT_EULA: "Y"
      MSSQL_PID: ${MSSQL_PID}
      MSSQL_SA_PASSWORD: ${MSSQL_SA_PASSWORD}
    ports: ["${MSSQL_PORT}:1433"]
    volumes: [mssql_data:/var/opt/mssql]

  orchestrator:                    # transporte Postgres <-> SQL Server
    build: ./orchestrator
    depends_on:
      postgres-oltp: { condition: service_healthy }
      postgres-dwh:  { condition: service_healthy }
      mssql-etl:     { condition: service_healthy }
    restart: unless-stopped

  data-gen:
    build: ./data-gen
    profiles: ["seed"]             # ne tourne que sur demande explicite
\end{lstlisting}

\paragraph{Healthchecks.} L'\code{orchestrator} ne démarre que lorsque
les trois bases passent leur \code{healthcheck} (\code{pg\_isready} pour
Postgres, \code{SELECT 1} via \code{sqlcmd} pour SQL Server). Le service
\code{data-gen}
utilise un \code{profile} pour ne pas se lancer automatiquement
(\code{docker compose --profile seed run --rm data-gen}).

\paragraph{Volumes nommés.} \code{oltp\_data}, \code{dwh\_data},
\code{mssql\_data}, \code{pgadmin\_data} : la persistance survit à
\code{docker compose down} mais peut être effacée par \code{down -v}.

\subsection{Configuration et démarrage}

Toute la configuration est centralisée dans un fichier \file{.env} versionné
qui contient les volumétries, les ports, les credentials SQL Server et les
paramètres cron de l'ETL.

\begin{lstlisting}[style=base,caption={Démarrage de la stack},label={lst:start}]
# 1) Bases + SQL Server + orchestrator + pgAdmin
docker compose up -d --build

# 2) Peuplement OLTP (one-shot)
docker compose --profile seed run --rm data-gen

# 3) Forcer un cycle ETL immediat (sinon il s'execute selon le cron)
docker compose restart orchestrator

# 4) Voir les logs ETL (orchestrateur Python + SQL Server)
docker compose logs -f orchestrator
docker compose logs -f mssql-etl
\end{lstlisting}

\screenshot{docker-startup.png}{Suivi du démarrage : \texttt{docker compose up -d --build} puis \texttt{docker compose logs -f orchestrator} permettent de vérifier l'enchaînement extract / EXEC sp\_run\_pipeline / push.}

\newpage
% =============================================================================
\section{Génération des données opérationnelles}
\label{sec:gen}
% =============================================================================

\subsection{Stratégie}

L'énoncé évoque 5\,500 produits, 90\,000 clients et 980\,000 commandes.
Pour rester dans des temps de cycle raisonnables en TP, le générateur
produit des volumes \emph{paramétrés} via le \file{.env} :

\begin{table}[h]
\centering\small
\begin{tabular}{lrr}\toprule
\textbf{Entité}    & \textbf{Cible énoncé} & \textbf{Défaut .env}\\\midrule
employes           &      $\approx$700      &       700\\
fournisseurs       &           64           &        64\\
produits           &      $\approx$5\,500    &     5\,500\\
clients            &      $\approx$90\,000   &    90\,000\\
commandes          &      $\approx$980\,000  &   980\,000\\
lignes\_commande   &      $\approx$2,5M     &  $\sim$2,5M\\
\bottomrule
\end{tabular}
\caption{Volumétries du générateur par défaut, alignées sur l'énoncé. Pour un cycle de TP plus rapide, ramener \texttt{GEN\_NB\_ORDERS} à 50\,000 dans \texttt{.env}.}
\end{table}

\subsection{Réalisme injecté}

\begin{itemize}
  \item Hiérarchie produit : 5 lignes $\times$ 12 catégories $\times$ 24 groupes,
    \emph{insérée par les seeds statiques}, puis affectation aléatoire des
    produits aux groupes.
  \item Historique de prix : 1 à 3 fenêtres temporelles par produit, marge
    cible entre 30 \% et 150 \% sur le coût.
  \item 30 \% des produits ont au moins une remise (5 à 30 \% sur 7 à 60 jours).
  \item 25 \% des clients possèdent une carte \emph{Orion Star Club}.
  \item Saisonnalité : surcharge de 30 \% sur les commandes des mois de
    novembre et décembre.
  \item Les commandes sont passées prioritairement par des employés du
    département \og Sales \fg{}.
  \item \code{total\_amount} est recalculé en SQL après l'insertion des lignes
    pour rester cohérent avec les remises appliquées.
\end{itemize}

\subsection{Listing : extrait du générateur}

\begin{lstlisting}[style=py,caption={Génération des prix historisés},label={lst:gen-prices}]
for pid in produit_ids:
    windows = []
    cursor_d = DATE_DEBUT - timedelta(days=180)
    while cursor_d < DATE_FIN:
        length = timedelta(days=random.randint(180, 720))
        end_d  = min(cursor_d + length, DATE_FIN + timedelta(days=365))
        windows.append((cursor_d, end_d))
        cursor_d = end_d + timedelta(days=1)
    for i, (sd, ed) in enumerate(windows):
        cout = round(random.uniform(5, 200), 2)
        prix = round(cout * random.uniform(1.3, 2.5), 2)
        # La derniere fenetre reste ouverte (date_fin IS NULL)
        end_val = None if i == len(windows) - 1 else ed
        prix_rows.append((pid, sd, end_val, cout, prix))
\end{lstlisting}

\newpage
% =============================================================================
\section{ETL T-SQL sur SQL Server}
\label{sec:etl}
% =============================================================================

\subsection{Architecture en deux niveaux}

L'ETL Orion est volontairement scindé en deux niveaux complémentaires :

\begin{table}[h]
\centering\small
\begin{tabularx}{\textwidth}{l X l}\toprule
\textbf{Couche}                  & \textbf{Rôle}                                            & \textbf{Techno}\\\midrule
Logique de transformation        & TRUNCATE/INSERT/MERGE, calcul SCD2, calcul des mesures, journal & \textbf{T-SQL} (procédures stockées sur SQL Server) \\
Transport / orchestration        & Extract Postgres OLTP, push staging SQL Server, EXEC procs, pull SQL Server, COPY DWH & Python (\code{orchestrate.py} + APScheduler)\\\bottomrule
\end{tabularx}
\caption{Séparation logique métier (T-SQL) vs transport (Python).}
\end{table}

Toute la \textbf{logique métier ETL} est concentrée dans des procédures
stockées T-SQL --- elles sont versionnées, testables individuellement
(\texttt{EXEC etl.sp\_load\_dim\_client}), et journalisées dans
\code{etl.run\_log}. L'orchestrateur Python est un \textbf{transporteur}
neutre, sans logique de transformation.

\begin{figure}[h]
\centering
\diag{07-etl-flow}
\caption{Pipeline ETL en 4 étapes : 1.~extract Postgres OLTP, 2.~staging SQL Server, 3.~transformation T-SQL (\texttt{etl.sp\_run\_pipeline}), 4.~load Postgres DWH.}
\end{figure}

\screenshot{ssms-objects.png}{SSMS --- arborescence \texttt{OrionETL} : schémas \texttt{staging}, \texttt{etl}, \texttt{dim}, \texttt{fact} avec leurs procédures stockées.}

\subsection{Procédures stockées T-SQL}

\begin{table}[h]
\centering\small
\begin{tabular}{l l l}\toprule
\textbf{Procédure}                  & \textbf{Mode}                   & \textbf{Cadence}\\\midrule
\code{etl.sp\_load\_dim\_date}      & bootstrap (idempotent)          & une fois\\
\code{etl.sp\_load\_dim\_canal}     & full reload                     & quotidien\\
\code{etl.sp\_load\_dim\_geographie}& full reload                     & quotidien\\
\code{etl.sp\_load\_dim\_fournisseur}& full reload                    & quotidien\\
\code{etl.sp\_load\_dim\_produit}   & full reload (SCD1)              & quotidien\\
\code{etl.sp\_load\_dim\_client}    & \textbf{SCD2 merge} (hash\_ligne)& quotidien\\
\code{etl.sp\_load\_dim\_employe}   & \textbf{SCD2 merge} (hash\_ligne)& quotidien\\
\code{etl.sp\_load\_fait\_ventes}   & \textbf{incrémental} (watermark)& quotidien\\
\code{etl.sp\_run\_pipeline}        & enchaîne tout                   & quotidien\\\bottomrule
\end{tabular}
\caption{Procédures stockées T-SQL et fréquence d'exécution.}
\end{table}

\subsection{Cœur SCD2 en T-SQL}

\begin{lstlisting}[style=sql,caption={Mise à jour SCD2 dans \texttt{sp\_load\_dim\_client} (extrait)},label={lst:scd2}]
;WITH src AS (
  SELECT client_id, nom_complet, sexe, tranche_age, groupe_client, ...,
         CONVERT(CHAR(64), HASHBYTES('SHA2_256',
           CONCAT_WS(N'|', nom_complet, sexe, tranche_age, groupe_client, ...)
         ), 2) AS hash_ligne
  FROM   staging.client_full
)
-- 1) Fermer la version courante quand le hash change
UPDATE d
   SET effectif_au = @ajd, est_courant = 0
  FROM dim.dim_client d
  JOIN src            s ON s.client_id = d.client_id
 WHERE d.est_courant = 1 AND d.hash_ligne <> s.hash_ligne;

-- 2) Inserer la nouvelle version (ou la premiere version)
INSERT dim.dim_client (..., effectif_du, effectif_au, est_courant, hash_ligne)
SELECT s.client_id, ..., @ajd, NULL, 1, s.hash_ligne
  FROM src s
  LEFT JOIN dim.dim_client d
    ON d.client_id = s.client_id AND d.est_courant = 1
 WHERE d.cle_client IS NULL OR d.hash_ligne <> s.hash_ligne;
\end{lstlisting}

\screenshot{ssms-sp-execute.png}{Exécution manuelle d'\texttt{etl.sp\_run\_pipeline} dans SSMS : sortie des PRINT et durée par procédure.}

\subsection{Watermark et idempotence}

\paragraph{Principe.} La table \code{etl.watermark(job\_name, last\_value)}
côté SQL Server mémorise pour chaque job la \emph{borne haute} de la dernière
extraction. Le job \code{sp\_load\_fact\_sales} consomme uniquement les lignes
\code{order\_date > watermark}, puis met à jour le watermark via
\code{MERGE} à \code{MAX(order\_date)} ingéré. Cela garantit :

\begin{itemize}
  \item \textbf{Pas de doublon} : la procédure commence par un
    \code{DELETE} ciblé sur \code{(order\_id, line\_no)} en cas de re-jeu
    accidentel d'un même lot.
  \item \textbf{Pas de perte} : c'est \code{order\_date} qui est tracké car
    il représente la dimension métier ; les commandes \og back-datées \fg{}
    sont rares dans Orion (énoncé).
  \item \textbf{Reprise sur erreur} : tout est dans une transaction T-SQL ;
    en cas d'échec, le watermark n'est pas avancé.
\end{itemize}

\subsection{Orchestrateur Python}

L'orchestrateur (\file{orchestrator/orchestrate.py}, $\sim$200 lignes) fait
trois choses : \emph{extract}, \emph{exec}, \emph{push}. Il est planifié
par \textbf{APScheduler} (cron-like in-process) ; pas d'Airflow, pas de Cron
système.

\begin{lstlisting}[style=py,caption={Cycle ETL côté orchestrateur},label={lst:orchestrate}]
def run_full_pipeline():
    with both_connections() as (pg_oltp, pg_dwh, ms):
        load_staging(pg_oltp, ms)   # 1) extract OLTP -> staging SQL Server
        run_pipeline(ms)            # 2) EXEC etl.sp_run_pipeline (T-SQL)
        push_to_dwh(ms, pg_dwh)     # 3) export dim/fact -> Postgres DWH

sched = BlockingScheduler(timezone="UTC")
sched.add_job(run_full_pipeline,
              CronTrigger(hour=ETL_HOUR, minute=ETL_MINUTE),
              id="etl_daily")
\end{lstlisting}

\screenshot{orchestrator-logs.png}{\texttt{docker compose logs -f orchestrator} --- une exécution complète : nombre de lignes par staging, durée totale, lignes \texttt{COPY} dans le DWH.}

\subsection{Journalisation}

Chaque procédure ouvre une ligne dans \code{etl.run\_log} via
\code{etl.sp\_run\_start} et la clôt via \code{etl.sp\_run\_end} avec les
champs \code{rows\_in}, \code{rows\_out}, \code{error\_msg} :

\begin{lstlisting}[style=sql,caption={Journal d'exécution côté SQL Server}]
SELECT TOP 20 job_name, started_at, ended_at, status,
       rows_in, rows_out, LEFT(error_msg,80) AS error_msg
FROM   OrionETL.etl.run_log
ORDER  BY run_id DESC;
\end{lstlisting}

\screenshot{ssms-run-log.png}{Tableau de bord ETL --- résultat de la requête sur \texttt{etl.run\_log} dans SSMS, état des derniers runs.}

\subsection{Variante \og tout SQL Server \fg{}}

L'image \code{mssql-etl} embarque déjà \code{unixODBC} et le driver
\code{psqlodbc} : on peut donc créer un \textbf{Linked Server} PostgreSQL et
faire des \code{OPENQUERY('ORION\_OLTP\_LINK', '...')} directement depuis
T-SQL, supprimant l'orchestrateur Python. Cette variante n'a pas été retenue
pour le TP (configuration ODBC plus fragile, moins de visibilité dans les
logs Docker), mais le rapport documente l'option.

\newpage
% =============================================================================
\section{Réponses aux questions analytiques}
\label{sec:answers}
% =============================================================================

Les requêtes ci-dessous tournent sur \code{orion\_dwh} (schéma \code{dw}).
Elles sont également regroupées dans \file{analytics/queries.sql} pour
exécution en lot.

% --- Q1 ----------------------------------------------------------------------
\question{Quels sont les produits qui se vendent le mieux ?}
\textit{Top 20 par quantité vendue, toutes périodes confondues.}
\begin{lstlisting}[style=sql]
SELECT  p.product_id, p.product_name, p.product_group_name,
        SUM(f.quantity)   AS qty_total,
        SUM(f.net_amount) AS revenue
FROM    fact_sales f JOIN dim_product p USING (product_key)
GROUP BY p.product_id, p.product_name, p.product_group_name
ORDER BY qty_total DESC
LIMIT 20;
\end{lstlisting}
\textbf{Lecture.} \code{qty\_total} mesure le volume, \code{revenue} la valeur.
On peut classer alternativement par \code{revenue} pour répondre \og en
valeur \fg{} plutôt qu'en \og en volume \fg{}.

% --- Q2 ----------------------------------------------------------------------
\question{Quels sont les produits en perte de vitesse ?}
\textit{Variation négative de CA d'une année sur l'autre, pire descente en tête.}
\begin{lstlisting}[style=sql]
WITH yearly AS (
  SELECT p.product_id, p.product_name, d.year_num,
         SUM(f.net_amount) AS revenue
  FROM   fact_sales f
  JOIN   dim_product p USING (product_key)
  JOIN   dim_date    d USING (date_key)
  GROUP  BY p.product_id, p.product_name, d.year_num
), delta AS (
  SELECT product_id, product_name, year_num, revenue,
         LAG(revenue) OVER (PARTITION BY product_id
                            ORDER BY year_num) AS revenue_prev
  FROM yearly
)
SELECT product_id, product_name, year_num,
       revenue, revenue_prev,
       revenue - revenue_prev AS delta_abs,
       CASE WHEN revenue_prev > 0
            THEN ROUND(100*(revenue - revenue_prev)/revenue_prev, 2)
       END AS delta_pct
FROM   delta
WHERE  revenue_prev IS NOT NULL AND revenue < revenue_prev
ORDER  BY delta_pct ASC NULLS LAST
LIMIT 20;
\end{lstlisting}

% --- Q3 ----------------------------------------------------------------------
\question{Produits qui contribuent très peu au CA pour un pays / une année donnés ?}
\textit{Critère : moins de 0,1 \% du CA pays-année.}
\begin{lstlisting}[style=sql]
WITH base AS (
  SELECT p.product_id, p.product_name, SUM(f.net_amount) AS revenue
  FROM   fact_sales f
  JOIN   dim_product   p USING (product_key)
  JOIN   dim_geography g ON g.geography_key = f.geography_cust_key
  JOIN   dim_date      d USING (date_key)
  WHERE  g.country_name = 'France'   -- :pays
    AND  d.year_num     = 2002       -- :annee
  GROUP  BY p.product_id, p.product_name
), total AS (SELECT SUM(revenue) AS tot FROM base)
SELECT b.product_id, b.product_name, b.revenue,
       ROUND(100*b.revenue/NULLIF(t.tot,0), 4) AS pct_ca
FROM   base b CROSS JOIN total t
WHERE  b.revenue/NULLIF(t.tot,0) < 0.001
ORDER  BY pct_ca ASC;
\end{lstlisting}
\textbf{Décision commerciale.} Cette même requête sert à proposer une remise
saisonnière sur les produits identifiés (insertion automatique dans
\code{ops.product\_discount} via un job dédié, hors périmètre du TP).

% --- Q4 ----------------------------------------------------------------------
\question{Marge générée par un groupe de produits, par année}
\begin{lstlisting}[style=sql]
SELECT  p.product_group_name, d.year_num,
        SUM(f.net_amount)    AS revenue,
        SUM(f.cost_amount)   AS cogs,
        SUM(f.margin_amount) AS margin,
        ROUND(100*SUM(f.margin_amount)/NULLIF(SUM(f.net_amount),0), 2) AS margin_pct
FROM    fact_sales f
JOIN    dim_product p USING (product_key)
JOIN    dim_date    d USING (date_key)
WHERE   p.product_group_name = 'Premium Hiking'    -- :groupe
GROUP BY p.product_group_name, d.year_num
ORDER BY d.year_num;
\end{lstlisting}

% --- Q5 ----------------------------------------------------------------------
\question{La marge dépend-elle de la quantité vendue ?}
\textit{Corrélation de Pearson + régression linéaire avec les fonctions Postgres natives.}
\begin{lstlisting}[style=sql]
WITH per_product AS (
  SELECT p.product_id,
         SUM(f.quantity)      AS qty,
         SUM(f.margin_amount) AS margin
  FROM   fact_sales f JOIN dim_product p USING (product_key)
  GROUP  BY p.product_id
)
SELECT  CORR(qty, margin)             AS pearson_qty_margin,
        REGR_SLOPE(margin, qty)       AS slope,
        REGR_INTERCEPT(margin, qty)   AS intercept,
        REGR_R2(margin, qty)          AS r_squared
FROM    per_product;
\end{lstlisting}

% --- Q6/Q7 -------------------------------------------------------------------
\question{Les remises font-elles augmenter les ventes (Q6) et la marge (Q7) ?}
\textit{Comparaison directe des lignes \emph{avec} et \emph{sans} remise.}
\begin{lstlisting}[style=sql]
SELECT  CASE WHEN discount_pct > 0 THEN 'avec_remise' ELSE 'sans_remise' END AS bucket,
        COUNT(*)                                AS nb_lignes,
        ROUND(AVG(quantity)::numeric,     3)    AS qty_moy_ligne,
        ROUND(AVG(net_amount)::numeric,   2)    AS ca_moy_ligne,
        ROUND(AVG(margin_amount)::numeric,2)    AS marge_moy_ligne,
        ROUND(SUM(net_amount)::numeric,   2)    AS ca_total,
        ROUND(SUM(margin_amount)::numeric,2)    AS marge_totale
FROM    fact_sales
GROUP   BY 1;
\end{lstlisting}
\textbf{Limite.} Cette comparaison est descriptive, non causale : pour
établir un effet causal il faudrait un dispositif d'A/B testing ou
un \emph{difference-in-differences} sur des produits pilotes.

% --- Q8 ----------------------------------------------------------------------
\question{Commerciaux qui font le plus de ventes}
\begin{lstlisting}[style=sql]
SELECT  e.employee_id, e.full_name, e.org_country,
        COUNT(DISTINCT f.order_id) AS nb_commandes,
        SUM(f.net_amount)          AS ca,
        SUM(f.margin_amount)       AS marge
FROM    fact_sales f
JOIN    dim_employee e ON e.employee_key = f.employee_key
WHERE   e.is_current
GROUP   BY 1, 2, 3
ORDER   BY ca DESC
LIMIT 20;
\end{lstlisting}

% --- Q9 ----------------------------------------------------------------------
\question{Commerciaux les plus performants par pays / sexe / âge / salaire}
\begin{lstlisting}[style=sql]
WITH base AS (
  SELECT e.employee_id, e.full_name, e.org_country,
         e.gender, e.age_band, e.salary_band,
         SUM(f.net_amount) AS ca
  FROM   fact_sales f JOIN dim_employee e ON e.employee_key = f.employee_key
  WHERE  e.is_current
  GROUP  BY 1, 2, 3, 4, 5, 6
)
SELECT org_country, gender, age_band, salary_band,
       employee_id, full_name, ca,
       ROW_NUMBER() OVER (PARTITION BY org_country  ORDER BY ca DESC) AS rang_pays,
       ROW_NUMBER() OVER (PARTITION BY gender       ORDER BY ca DESC) AS rang_sexe,
       ROW_NUMBER() OVER (PARTITION BY age_band     ORDER BY ca DESC) AS rang_age,
       ROW_NUMBER() OVER (PARTITION BY salary_band  ORDER BY ca DESC) AS rang_salaire
FROM   base
ORDER  BY ca DESC;
\end{lstlisting}

% --- Q10 ---------------------------------------------------------------------
\question{Quels groupes de clients sont identifiés ?}
\begin{lstlisting}[style=sql]
SELECT c.customer_group,
       COUNT(DISTINCT c.customer_id) AS nb_clients,
       SUM(f.net_amount)             AS ca,
       ROUND(AVG(f.net_amount)::numeric, 2) AS panier_moyen_ligne
FROM   fact_sales f JOIN dim_customer c ON c.customer_key = f.customer_key
WHERE  c.is_current
GROUP  BY c.customer_group
ORDER  BY ca DESC;
\end{lstlisting}
\textbf{Note.} Les groupes de référence sont \emph{Bronze}, \emph{Silver},
\emph{Gold}, \emph{VIP}, \emph{Inactif}, peuplés dans le fichier de seeds
statiques de l'OLTP.

% --- Q11 ---------------------------------------------------------------------
\question{Clients les plus rentables (top 50)}
\begin{lstlisting}[style=sql]
SELECT c.customer_id, c.full_name, c.country_name,
       SUM(f.net_amount)    AS ca,
       SUM(f.margin_amount) AS marge,
       COUNT(DISTINCT f.order_id) AS nb_commandes
FROM   fact_sales f JOIN dim_customer c ON c.customer_key = f.customer_key
WHERE  c.is_current
GROUP  BY 1, 2, 3
ORDER  BY marge DESC
LIMIT 50;
\end{lstlisting}

% --- Q12 ---------------------------------------------------------------------
\question{Fournisseurs proposant des produits rentables}
\begin{lstlisting}[style=sql]
SELECT s.supplier_id, s.supplier_name, s.country_name,
       SUM(f.net_amount)    AS ca,
       SUM(f.margin_amount) AS marge,
       ROUND(100*SUM(f.margin_amount)/NULLIF(SUM(f.net_amount),0),2) AS taux_marge
FROM   fact_sales f JOIN dim_supplier s USING (supplier_key)
GROUP  BY 1, 2, 3
ORDER  BY marge DESC
LIMIT 20;
\end{lstlisting}

% --- Q13 ---------------------------------------------------------------------
\question{Moyenne et écart-type du chiffre d'affaires}
\textit{Trois grains : par commande, par mois, par employé.}
\begin{lstlisting}[style=sql]
-- Par commande
WITH per_order AS (
  SELECT order_id, SUM(net_amount) AS ca FROM fact_sales GROUP BY order_id
)
SELECT AVG(ca) AS ca_moy, STDDEV_SAMP(ca) AS ca_std,
       MIN(ca) AS ca_min, MAX(ca) AS ca_max
FROM   per_order;

-- Par mois
WITH per_month AS (
  SELECT d.year_num, d.month_num, SUM(f.net_amount) AS ca
  FROM   fact_sales f JOIN dim_date d USING (date_key)
  GROUP  BY d.year_num, d.month_num
)
SELECT AVG(ca) AS ca_moy_mois, STDDEV_SAMP(ca) AS ca_std_mois FROM per_month;

-- Par employe
WITH per_emp AS (
  SELECT e.employee_id, SUM(f.net_amount) AS ca
  FROM   fact_sales f JOIN dim_employee e ON e.employee_key = f.employee_key
  WHERE  e.is_current
  GROUP  BY e.employee_id
)
SELECT AVG(ca) AS ca_moy_emp, STDDEV_SAMP(ca) AS ca_std_emp FROM per_emp;
\end{lstlisting}

% --- Q14 ---------------------------------------------------------------------
\question{Variables qui expliquent le mieux le chiffre d'affaires}
\textit{Calcul du $\eta^2$ (eta-squared) = part de variance inter-groupes.
Plus $\eta^2$ est proche de 1, plus la variable est explicative.}
\[
\eta^2 = \frac{\sum_g n_g (\bar y_g - \bar y)^2}{\sum_i (y_i - \bar y)^2}
\]
\begin{lstlisting}[style=sql]
WITH facts AS (
  SELECT f.net_amount,
         p.product_line_name, p.product_category_name,
         c.customer_group, ch.channel_code,
         e.gender AS emp_gender, e.salary_band AS emp_salary_band,
         g.country_name AS cust_country
  FROM   fact_sales   f
  JOIN   dim_product  p  USING (product_key)
  JOIN   dim_customer c  ON c.customer_key = f.customer_key
  JOIN   dim_channel  ch USING (channel_key)
  JOIN   dim_employee e  ON e.employee_key = f.employee_key
  JOIN   dim_geography g ON g.geography_key = f.geography_cust_key
), mu_t AS (SELECT AVG(net_amount) AS m_t FROM facts),
   sst  AS (SELECT SUM(power(net_amount-m_t,2)) AS st FROM facts, mu_t),
   eta(variable, eta_squared) AS (
     SELECT 'product_line', (
       SELECT SUM(n_g*power(m_g-m_t,2))/NULLIF(st,0)
       FROM (SELECT product_line_name, AVG(net_amount) m_g, COUNT(*) n_g
             FROM facts GROUP BY product_line_name) g, mu_t, sst)
     UNION ALL SELECT 'customer_group', (...)   -- etc.
   )
SELECT variable, ROUND(eta_squared::numeric, 4) AS eta_squared
FROM   eta ORDER BY eta_squared DESC NULLS LAST;
\end{lstlisting}
\textbf{Lecture.} On obtient un classement des variables candidates par
pouvoir explicatif. Pour aller plus loin (régression linéaire ou
non-linéaire), on exporte vers Python/R via Pandas.

% --- Q15 ---------------------------------------------------------------------
\question{Différence significative de CA entre commerciaux femmes et hommes ?}
\textit{Statistique de Welch (test $t$ pour variances inégales).}
\begin{lstlisting}[style=sql]
WITH per_emp AS (
  SELECT e.gender, e.employee_id, SUM(f.net_amount) AS ca
  FROM   fact_sales f JOIN dim_employee e ON e.employee_key = f.employee_key
  WHERE  e.is_current
  GROUP  BY e.gender, e.employee_id
), agg AS (
  SELECT gender, COUNT(*) AS n, AVG(ca) AS mean_ca, VAR_SAMP(ca) AS var_ca
  FROM   per_emp GROUP BY gender
)
SELECT f.mean_ca AS mean_F, m.mean_ca AS mean_M,
       f.n      AS n_F,    m.n      AS n_M,
       (f.mean_ca - m.mean_ca)
       / sqrt(f.var_ca/f.n + m.var_ca/m.n) AS welch_t
FROM   (SELECT * FROM agg WHERE gender='F') f
CROSS JOIN (SELECT * FROM agg WHERE gender='M') m;
\end{lstlisting}
\textbf{Décision.} La p-value associée à $t$ se calcule avec la loi de
Student aux degrés de liberté de Welch ; comme PostgreSQL n'embarque pas
\code{pt()}, on le fait en SciPy après export, ou on installe l'extension
\code{tablefunc} / \code{pgsl-stats}.

\newpage
% =============================================================================
\section*{Annexe A --- Arborescence du projet}
\addcontentsline{toc}{section}{Annexe A --- Arborescence du projet}
% =============================================================================

\begin{lstlisting}[style=base,basicstyle=\ttfamily\scriptsize]
tp1-orion/
|-- README.md
|-- docker-compose.yml
|-- .env
|-- doc/
|   |-- MODELISATION.md            (modele OP + transition + DW)
|   |-- UML.md                     (11 diagrammes UML)
|   |-- pgadmin-servers.json       (auto-config pgAdmin)
|   `-- rapport/                   (ce document LaTeX)
|       |-- rapport.tex            (rapport principal A4)
|       |-- rapport.pdf
|       |-- uml-poster.tex         (poster A0 : 11 diag. UML)
|       |-- uml-individual.tex     (1 page A3 par diagramme)
|       |-- diagrams/              (sources Mermaid .mmd)
|       |-- screenshots/           (PNG des captures d'ecran)
|       `-- build/                 (PDF compiles)
|-- oltp/
|   `-- init/
|       |-- 01_schema.sql          (DDL OLTP, schema "ops")
|       `-- 02_static_data.sql     (referentiels statiques)
|-- dwh/
|   `-- init/
|       `-- 01_schema.sql          (DDL DWH, schema "dw")
|-- data-gen/
|   |-- Dockerfile
|   |-- requirements.txt
|   `-- generate.py                (Faker, peuplement OLTP)
|-- mssql-etl/                     ** moteur ETL : SQL Server 2022 **
|   |-- Dockerfile                 (mssql/server + unixODBC + psqlodbc)
|   |-- scripts/
|   |   |-- entrypoint.sh          (start sqlservr, run init/*.sql)
|   |   `-- init-db.sh             (re-execution manuelle)
|   `-- init/                      (7 scripts T-SQL, executes en ordre)
|       |-- 01_database.sql        (CREATE DATABASE OrionETL + schemas)
|       |-- 02_meta_tables.sql     (etl.run_log, etl.watermark, helpers)
|       |-- 03_staging.sql         (staging.* miroir denormalise OLTP)
|       |-- 04_dim_tables.sql      (dim.* + fact.* gold)
|       |-- 05_sp_dim_simple.sql   (sp_load_dim_date/channel/geo/...)
|       |-- 06_sp_dim_scd2.sql     (sp_load_dim_customer/employee SCD2)
|       `-- 07_sp_fact.sql         (sp_load_fact_sales + sp_run_pipeline)
|-- orchestrator/                  ** transport Postgres <-> SQL Server **
|   |-- Dockerfile                 (python:3.12 + msodbcsql18)
|   |-- requirements.txt
|   `-- orchestrate.py             (extract / EXEC / push, APScheduler)
`-- analytics/
    `-- queries.sql                (15 requetes repondant aux questions)
\end{lstlisting}

% =============================================================================
\section*{Annexe B --- Démarrage et exploitation}
\addcontentsline{toc}{section}{Annexe B --- Démarrage et exploitation}
% =============================================================================

\begin{lstlisting}[style=base,caption={Cycle de vie complet}]
# Demarrer (build des images mssql-etl + orchestrator inclus)
cd ~/dev-laboratory/tp-ed-eneam/tp1-orion
docker compose up -d --build

# Peupler l'OLTP avec Faker (one-shot)
docker compose --profile seed run --rm data-gen

# Forcer un cycle ETL immediat
docker compose restart orchestrator

# Inspecter
docker compose ps
docker compose logs -f orchestrator
docker compose logs -f mssql-etl

# Journal SQL Server
docker compose exec mssql-etl /opt/mssql-tools18/bin/sqlcmd \
    -S localhost -U sa -P "$MSSQL_SA_PASSWORD" -No -d OrionETL \
    -Q "SELECT TOP 10 job_name, started_at, status, rows_in, rows_out
        FROM etl.run_log ORDER BY run_id DESC"

# Resultat dans le DWH
docker compose exec postgres-dwh psql -U dwh -d orion_dwh \
    -c "SELECT COUNT(*) FROM dw.fait_ventes;"

# Reset complet (efface les volumes)
docker compose down -v
\end{lstlisting}

\screenshot{pgadmin-dwh-fact.png}{pgAdmin --- aperçu du DWH après chargement : ligne d'\texttt{dw.fait\_ventes} avec ses surrogate keys et mesures.}

\screenshot{queries-result.png}{Exécution d'une requête analytique (Q1 : produits qui se vendent le mieux) sur \texttt{orion\_dwh}.}

\bigskip
\hrule
\bigskip
\centering
{\itshape\small Fin du rapport.}

\end{document}
